% ============================================================
%  ASPETTO DEI LISTATI: UNA FINESTRA DI VS CODE (TEMA DARK+)
% ============================================================
% Lo stesso aspetto della dispensa di Tecnologie Web: ogni listato è una
% piccola finestra dell'editor Visual Studio Code con il tema scuro
% predefinito (Dark+). Sfondo quasi nero, numeri di riga grigi a sinistra e,
% se il listato ha un nome (chiave "name"), la barra delle schede in alto con
% la scheda del file aperto: icona del linguaggio, nome del file e crocetta.
%
% Uso:  \begin{lstlisting}[style=java, name=Ordine.java] ... \end{lstlisting}
% Stili: java, json, http, shell, yaml, docker, conf, xml, js, c, plain.

\definecolor{vscbg}{HTML}{1E1E1E}        % sfondo dell'editor
\definecolor{vsctabbar}{HTML}{252526}    % barra delle schede
\definecolor{vscfg}{HTML}{D4D4D4}        % testo normale, punteggiatura
\definecolor{vsclinenum}{HTML}{858585}   % numeri di riga
\definecolor{vscblue}{HTML}{569CD6}      % class, public, int, true...
\definecolor{vscpurple}{HTML}{C586C0}    % if, return, new, throw... (controllo)
\definecolor{vscstring}{HTML}{CE9178}    % stringhe
\definecolor{vsccomment}{HTML}{6A9955}   % commenti
\definecolor{vscfunc}{HTML}{DCDCAA}      % nomi di metodo
\definecolor{vscvar}{HTML}{9CDCFE}       % variabili, campi, chiavi
\definecolor{vsctype}{HTML}{4EC9B0}      % classi, interfacce, annotazioni
\definecolor{vscnum}{HTML}{B5CEA8}       % numeri
\definecolor{vscaccent}{HTML}{0078D4}    % filetto della scheda attiva
\definecolor{vsctabtext}{HTML}{FFFFFF}   % nome del file nella scheda
\definecolor{vsctabdim}{HTML}{9D9D9D}    % crocetta di chiusura
% colori delle icone dei file (tema di icone Seti, quello predefinito)
\definecolor{vsciconorange}{HTML}{E37933}
\definecolor{vsciconblue}{HTML}{519ABA}
\definecolor{vsciconyellow}{HTML}{CBCB41}
\definecolor{vscicongreen}{HTML}{8DC149}
\definecolor{vsciconpurple}{HTML}{A074C4}
\definecolor{vsciconred}{HTML}{CC3E44}

% ============================================================
%  LINGUAGGI
% ============================================================
% Listings non conosce la grammatica: colora le parole che compaiono in
% elenchi fissi. Le classi di parole chiave riproducono le categorie di Dark+:
%   [1] azzurro   dichiarazioni, modificatori, tipi primitivi, letterali
%   [2] viola     controllo del flusso (if, return, new, throw, try...)
%   [3] giallo    metodi chiamati o dichiarati negli esempi della dispensa
%   [4] verde     classi, interfacce, tipi e annotazioni
% VS Code colora di giallo qualunque nome seguito da "(": listings non può
% guardare avanti, quindi l'elenco [3] contiene i nomi dei metodi usati negli
% esempi. Tutti gli altri identificatori sono variabili (azzurro chiaro).
\lstset{keywords={}}

% ─── Java ────────────────────────────────────────────────────
\lstdefinelanguage{JavaVSC}{
  sensitive=true,
  alsoletter={@},
  morecomment=[l]{//},
  morecomment=[s]{/*}{*/},
  morestring=[b]",
  morestring=[b]',
  keywords=[1]{
    abstract,boolean,byte,char,class,double,enum,extends,final,float,
    implements,import,int,interface,long,native,package,private,protected,
    public,record,short,static,synchronized,this,super,throws,transient,
    void,volatile,var,true,false,null,sealed,permits},
  keywords=[2]{
    return,if,else,for,while,do,switch,case,default,break,continue,try,
    catch,finally,throw,new,yield,assert,instanceof},
  % i metodi degli esempi sono raccolti da preamble/nomi_listati.py (vedi listati-nomi.tex)
  keywords=[3]{main,println,print,printf,format},
  keywords=[4]{
    String,Object,Integer,Long,Double,Boolean,Character,Math,System,
    List,ArrayList,LinkedList,Map,HashMap,Set,HashSet,Optional,Collection,
    Collections,Arrays,Iterator,Iterable,Comparable,Comparator,Runnable,
    Thread,Exception,RuntimeException,IllegalArgumentException,
    IllegalStateException,NullPointerException,IOException,
    UnsupportedOperationException,
    Queue,BlockingQueue,LinkedBlockingQueue,AtomicInteger,Instant,Duration,
    LocalDate,LocalDateTime,UUID,BigDecimal,Stream,Collectors,
    HttpClient,HttpRequest,HttpResponse,URI,Base64,Mac,SecretKeySpec,
    StandardCharsets,
    Override,Deprecated,FunctionalInterface,Test,BeforeEach,
    RestController,RequestMapping,GetMapping,PostMapping,PutMapping,
    DeleteMapping,PathVariable,RequestBody,RequestParam,RequestHeader,
    ResponseEntity,Service,Repository,Component,Autowired,Entity,Id,
    GeneratedValue,SpringBootApplication,SpringApplication,
    Assertions},
}

% ─── C ───────────────────────────────────────────────────────
\lstdefinelanguage{CVSC}{
  sensitive=true,
  morecomment=[l]{//},
  morecomment=[s]{/*}{*/},
  morecomment=[l][\color{vscpurple}]{\#},
  morestring=[b]",
  morestring=[b]',
  keywords=[1]{int,char,float,double,void,long,short,unsigned,signed,
    struct,union,enum,typedef,const,static,extern,sizeof,NULL},
  keywords=[2]{return,if,else,for,while,do,switch,case,default,break,
    continue,goto},
  keywords=[3]{printf,scanf,malloc,free,strcpy,strlen,memcpy,main},
}

% ─── JavaScript (lato client delle API REST) ─────────────────
\lstdefinelanguage{JSVSC}{
  sensitive=true,
  morecomment=[l]{//},
  morecomment=[s]{/*}{*/},
  morestring=[b]",
  morestring=[b]',
  morestring=[b]`,
  keywords=[1]{var,let,const,function,class,new,this,async,null,undefined,
    true,false},
  keywords=[2]{return,if,else,for,while,try,catch,throw,await,import,
    export,from,of},
  keywords=[3]{fetch,json,log,then,stringify,parse,getItem,setItem},
  keywords=[4]{console,JSON,localStorage,Error},
}

% ─── JSON ────────────────────────────────────────────────────
\lstdefinelanguage{JSONVSC}{
  sensitive=true,
  morestring=[b]",
  morekeywords={true,false,null},
}

% ─── XML ─────────────────────────────────────────────────────
\lstloadaspects{html}
\lstdefinelanguage{XMLVSC}{
  sensitive=true,
  morestring=[b]",
  tag=**[s]<>,
  usekeywordsintag=false,
  markfirstintag=true,
  morecomment=[s]{<!--}{-->},
}

% ─── HTTP (messaggi grezzi di richiesta e risposta) ──────────
\lstdefinelanguage{HTTPVSC}{
  sensitive=true,
  alsoletter={/.-},
  morestring=[b]",
  keywords=[1]{GET,POST,PUT,PATCH,DELETE,HEAD,OPTIONS},
  keywords=[2]{HTTP/1.0,HTTP/1.1,HTTP/2},
  keywords=[3]{Host,User-Agent,Accept,Content-Type,Content-Length,
    Location,Authorization,WWW-Authenticate,Cache-Control,Date,Server,
    Allow,ETag,Connection},
}

% ─── Terminale ───────────────────────────────────────────────
% I comandi sono gialli, come le funzioni (VS Code fa lo stesso).
\lstdefinelanguage{ShellVSC}{
  sensitive=true,
  morecomment=[l]{\#},
  morestring=[b]",
  morestring=[b]',
  morekeywords={docker,curl,aws,az,gcloud,kubectl,java,javac,mvn,git,
    ssh,cd,ls,cat,echo,export,sudo,virsh,VBoxManage,qemu-img,terraform,
    base64,jq,openssl,npm,node},
}

% ─── YAML (docker compose, Kubernetes, configurazioni) ───────
\lstdefinelanguage{YAMLVSC}{
  sensitive=true,
  morecomment=[l]{\#},
  morestring=[b]",
  morestring=[b]',
  morekeywords={true,false,null},
}

% ─── File di configurazione (nginx e simili) ─────────────────
\lstdefinelanguage{ConfVSC}{
  sensitive=true,
  morecomment=[l]{\#},
  morestring=[b]",
  morekeywords={upstream,server,listen,location,proxy_pass,least_conn,backup},
}

% ─── Dockerfile ──────────────────────────────────────────────
\lstdefinelanguage{DockerVSC}{
  sensitive=false,
  morecomment=[l]{\#},
  morestring=[b]",
  morekeywords={FROM,AS,RUN,CMD,COPY,ADD,WORKDIR,EXPOSE,ENV,ARG,
    ENTRYPOINT,USER,VOLUME,LABEL,HEALTHCHECK},
}

% I nomi di metodi e tipi usati negli esempi, raccolti da preamble/nomi_listati.py
% File generato da preamble/nomi_listati.py: non modificarlo a mano.
% Metodi (giallo) e tipi (verde acqua) che compaiono nei listati Java della dispensa.
\lstdefinelanguage{JavaVSCnomi}[]{JavaVSC}{
  morekeywords=[3]{
    accept,acquisisciAssegno,actionPerformed,add,addActionListener,addebita,
    addebitaIlTotaleDelCarrello,addMoney,addSocieta,addSocietaDB,aggiungi,aggiungiSocieta,
    aggiungiSocietaBorsa,aggiungiStudente,allOf,annota,applicaLoScontoRigaPerRiga,
    applicaScontiInVigore,assertEquals,assertFalse,assertThrows,assertTrue,availableProcessors,
    avvia,azioneEvento1,azzera,balanceOn,battito,bisestile,body,build,calcola,calcolaTotaleCarrello,
    carica,caricaUtentiBloccati,carrelloDa,ceil,charAt,checkMirrors,chiudi,ciclo,comeRecord,
    comparing,compile,compute,conferma,connetti,contains,controlla,converti,create,created,
    createStatement,currentThread,currentTimeMillis,daJson,decode,delete,deleteById,dettaglio,
    dispose,doFinal,dopoOperazione,drive,elabora,email,emergenza,encode,encodeToString,equals,
    erogaContanti,esegui,executeQuery,executeUpdate,exit,fattureInSospeso,filter,find,findAll,
    findById,findByListing,firma,firstValue,flick,floorMod,format,formatReport,generaFatturaPdf,
    gestisci,get,getAndIncrement,getBytes,getCarburator,getConnection,getCredito,getData,getEngine,
    getEpochSecond,getFirstName,getInputStream,getInstance,getLastName,getListino,getMessage,
    getNome,getOutputStream,getParent,getPayment,getPrezzo,getPrezzoUnitario,getProperty,
    getQuantita,getRuntime,getSconto,getSpeedLimit,getText,getTotalMoney,getUrlDecoder,
    getUrlEncoder,getWallet,group,habitat,haPrenotato,hashCode,headers,html,identificaCarta,
    ilMilleNovecentoNonEraBisestile,increaseSpeed,info,init,inserisciMoneta,interrupt,invia,
    inviaEmail,inviaFattura,isAfter,isAperto,isDirectory,isEqual,issueToken,istanzeDesiderate,join,
    leggiQuotazione,list,listFiles,load,loSlotNonEPiuDisponibileAlMomentoDellaConferma,main,map,
    matcher,mostra,move,nanoTime,newBuilder,newFixedThreadPool,newHttpClient,noContent,
    nonCiSiPuoPrenotareDueVolte,notFound,notifica,now,nuovaSocieta,of,offer,ofMillis,ofSeconds,
    ofString,openFuelValve,orElse,orElseThrow,paga,pagata,parallel,parseInt,parseLong,pianifica,
    plus,preleva,prelevaProdotto,prelievo,prenota,prenotazioniAttive,prepara,prepareStatement,print,
    printf,println,prioritario,prodotto,quotazione,rangeClosed,readLine,registraVendita,renameTo,
    replace,resolve,reviewsOf,riceviPagamento,richiestaRestituzione,richiestaRestituzioneCredito,
    rifiuta,runAsync,saturaOSegnala,save,scegli,sceltaProdotto,search,segnaleRicevuto,
    selezionaProdotto,send,setCredito,setDouble,setState,setString,setTotalMoney,setV,setVisible,
    setWallet,shardPer,size,slotLiberi,sorted,spedizioneEspressa,spedizioneNormale,
    spegniDistributore,split,stampa,startJourney,statusCode,stream,submit,subtractMoney,sum,take,
    tempoMedioDiRispostaEntroCentoMillisecondi,timeout,timeOutErogazione,timeoutErogazione,toList,
    toUpperCase,trova,ultimaQuotazioneNota,ultimoImporto,ultimoNoto,verifyToken,versaAssegno,
    vuoleLaRicevuta,warn,withdraw,withoutPadding},
  morekeywords=[4]{
    @DeleteMapping,@GetMapping,@Override,@PathVariable,@PostMapping,@PutMapping,@RequestBody,
    @RequestMapping,@RequestParam,@RestController,@Test,Account,AccountingDriver,ActionEvent,
    ActionListener,AggiungiNuovaSocieta,Animal,Appello,AquaticAnimal,ArrayList,Assegno,Assertions,
    AtomicInteger,Attesa,AttesaSceltaProdotto,Automobile,Autoscaler,Badge,Bancomat,Base,Base64,
    BigDecimal,Bike,BlockingQueue,BodyHandlers,BodyPublishers,Borsa,BufferedReader,Calendario,Car,
    Carrello,CarrelloTest,Carta,Casa,Cashier,Cassa,Character,Cliente,Comparator,CompletableFuture,
    ConnessioneDatabase,Contabilita,Controller,ControlloMarcia,Corso,Counter,Customer,DataAccess,
    DataSource,Date,Dieta,Distributore,Double,Driver,DriverManager,Duration,Employee,Encoder,Entry,
    ErogazioneProdotto,Esame,Esito,EstensionePrelievo,Exception,Executors,ExecutorService,Fattura,
    File,FileInputStream,Files,GatewayFinto,GatewayPagamenti,HashSet,HttpClient,HttpRequest,
    HttpResponse,IllegalArgumentException,IllegalStateException,InputStreamReader,Instant,Integer,
    InterruptedException,IOException,Jwt,LandAnimal,LinkedBlockingQueue,List,ListinoDAO,
    ListinoImplementazionePostgresDAO,LocalTime,Long,LongStream,Mac,Machine,Marca,Matcher,Math,
    MessageDigest,Money,Nodo,OKButton,Optional,Ordine,OrdineController,OrdineRepository,
    OrdineService,Path,Pattern,Person,Persona,Prenotazione,PreparedStatement,PrintWriter,Prodotto,
    Properties,Quotazione,Report,ReportDataElement,ReportFormatter,ReportPrinter,ResponseEntity,
    ResultSet,Review,ReviewController,Ricevuta,RigaCarrello,Rinomina,Rinominatore,Road,RoundRobin,
    Route,Runnable,Runtime,Ruota,SecretKeySpec,Sentinella,ServerSocket,Set,Short,Slot,Societa,
    Socket,Somma,SQLException,Stampabile,StampaRicevuta,StandardCharsets,Stanza,State,Stream,String,
    Studente,System,Test,Thread,Todo,TodoClient,TodoController,TodoRepository,TransactionProcessor,
    Traveler,Vehicle,Veicolo,Vendita,VistaOrdine,Void,Voto,Wallet},
}


% ============================================================
%  STILI DI LISTATO
% ============================================================
% Ogni stile sceglie anche l'icona mostrata nella scheda (chiave vscicon).
\lstdefinestyle{java}{
  language=JavaVSCnomi,
  vscicon=java,
  keywordstyle=[1]\color{vscblue},
  keywordstyle=[2]\color{vscpurple},
  keywordstyle=[3]\color{vscfunc},
  keywordstyle=[4]\color{vsctype},
  identifierstyle=\color{vscvar},
  stringstyle=\color{vscstring},
  commentstyle=\color{vsccomment},
}
\lstdefinestyle{c}{
  language=CVSC,
  vscicon=c,
  keywordstyle=[1]\color{vscblue},
  keywordstyle=[2]\color{vscpurple},
  keywordstyle=[3]\color{vscfunc},
  identifierstyle=\color{vscvar},
  stringstyle=\color{vscstring},
  commentstyle=\color{vsccomment},
}
\lstdefinestyle{js}{
  language=JSVSC,
  vscicon=js,
  keywordstyle=[1]\color{vscblue},
  keywordstyle=[2]\color{vscpurple},
  keywordstyle=[3]\color{vscfunc},
  keywordstyle=[4]\color{vsctype},
  identifierstyle=\color{vscvar},
  stringstyle=\color{vscstring},
  commentstyle=\color{vsccomment},
}
% Nel JSON le chiavi sono azzurre e i valori arancioni: listings non le
% distingue (sono stringhe entrambe), quindi le stringhe restano arancioni.
\lstdefinestyle{json}{
  language=JSONVSC,
  vscicon=json,
  keywordstyle=\color{vscblue},
  stringstyle=\color{vscstring},
}
\lstdefinestyle{xml}{
  language=XMLVSC,
  vscicon=xml,
  keywordstyle=\color{vscblue},
  tagstyle=\color{vscvar},
  stringstyle=\color{vscstring},
  commentstyle=\color{vsccomment},
}
\lstdefinestyle{http}{
  language=HTTPVSC,
  vscicon=http,
  keywordstyle=[1]\color{vscpurple},
  keywordstyle=[2]\color{vscblue},
  keywordstyle=[3]\color{vscvar},
  stringstyle=\color{vscstring},
}
\lstdefinestyle{shell}{
  language=ShellVSC,
  vscicon=shell,
  keywordstyle=\color{vscfunc},
  stringstyle=\color{vscstring},
  commentstyle=\color{vsccomment},
}
\lstdefinestyle{yaml}{
  language=YAMLVSC,
  vscicon=yaml,
  keywordstyle=\color{vscblue},
  stringstyle=\color{vscstring},
  commentstyle=\color{vsccomment},
}
\lstdefinestyle{docker}{
  language=DockerVSC,
  vscicon=docker,
  keywordstyle=\color{vscpurple},
  stringstyle=\color{vscstring},
  commentstyle=\color{vsccomment},
}
\lstdefinestyle{conf}{
  language=ConfVSC,
  vscicon=conf,
  keywordstyle=\color{vscblue},
  stringstyle=\color{vscstring},
  commentstyle=\color{vsccomment},
}
\lstdefinestyle{plain}{
  language={},
  vscicon=plain,
}

% ============================================================
%  LA FINESTRA DELL'EDITOR
% ============================================================
\makeatletter
% Chiave vscicon: l'icona da mettere nella scheda (la sceglie lo stile)
\lst@Key{vscicon}{plain}{\def\igvsc@icon{#1}}

\newcommand{\igvsccodefont}{\fontsize{7.5}{9}\selectfont\FiraCodeNL\color{vscfg}}
\lstset{
  % la classe attiva mathescape: il "$" del prompt aprirebbe una formula
  mathescape=false,
  basicstyle=\linespread{1.15}\igvsccodefont,
  % literate={} annulla le sostituzioni dello stile FiraCodeStyle (pacchetto
  % lstfiracode, caricato dalla classe): con quelle listings non riconosce i
  % delimitatori che contengono le coppie di caratteri delle legature
  literate={},
  backgroundcolor=\color{vscbg},
  identifierstyle={},
  % i listati scritti con language=... e senza stile restano leggibili sullo
  % sfondo scuro
  keywordstyle=\color{vscblue},
  stringstyle=\color{vscstring},
  commentstyle=\color{vsccomment},
  % numeri di riga, come nella "gutter" di VS Code
  numbers=left,
  numberstyle=\fontsize{5.5}{6}\selectfont\FiraCodeNL\color{vsclinenum},
  numbersep=5pt,
  firstnumber=1,
  % righe lunghe: si va a capo con un rientro, invece di uscire dal margine
  breaklines=true,
  breakindent=1.5em,
  postbreak={},
  columns=fullflexible,
  keepspaces=true,
  showstringspaces=false,
  tabsize=2,
  % la cornice è invisibile: serve solo a estendere lo sfondo scuro attorno
  % al codice e sotto i numeri di riga
  frame=single,
  framerule=0pt,
  rulecolor=\color{vscbg},
  % il "titolo" del listato diventa la barra delle schede, sopra il codice
  captionpos=t,
  abovecaptionskip=0pt,
  belowcaptionskip=0pt,
  aboveskip=6pt,
  belowskip=6pt,
  % dentro la finestra: 16pt a sinistra per i numeri di riga, 4pt a destra
  xleftmargin=16pt, framexleftmargin=12pt, framesep=4pt,
  xrightmargin=4pt, framexrightmargin=0pt,
}

% Icone dei file nella scheda (riprendono il tema di icone Seti di VS Code)
\newcommand{\igvsc@badge}[2]{%
  {\FiraCodeNL\bfseries\fontsize{6.5}{7}\selectfont\color{#1}#2}}
\newcommand{\igvsc@iconmark}{%
  \ifcsname igvsc@icon@\igvsc@icon\endcsname
    \csname igvsc@icon@\igvsc@icon\endcsname
  \else
    \igvsc@icon@plain
  \fi}
\def\igvsc@icon@java{\igvsc@badge{vsciconred}{J}}
\def\igvsc@icon@c{\igvsc@badge{vsciconblue}{C}}
\def\igvsc@icon@js{\igvsc@badge{vsciconyellow}{JS}}
\def\igvsc@icon@json{\igvsc@badge{vsciconyellow}{\{\}}}
\def\igvsc@icon@xml{\igvsc@badge{vsciconorange}{<>}}
\def\igvsc@icon@shell{\igvsc@badge{vscicongreen}{>\_}}
\def\igvsc@icon@yaml{\igvsc@badge{vsciconpurple}{!}}
\def\igvsc@icon@docker{\igvsc@badge{vsciconblue}{D}}
\def\igvsc@icon@http{{\fontsize{7}{7}\selectfont\color{vsciconblue}\ifilecode}}
\def\igvsc@icon@conf{\igvsc@badge{vscicongreen}{\{\}}}
\def\igvsc@icon@plain{{\fontsize{7}{7}\selectfont\color{vsctabdim}\ifilecode}}

% La barra delle schede: grigio scuro con gli angoli superiori arrotondati, e
% a sinistra la scheda attiva (dello stesso colore dell'editor, con il
% filetto blu in alto) che contiene icona, nome del file e crocetta.
\newlength{\igvsc@barw}
\newcommand{\igvsc@tabbar}{%
  \setlength{\igvsc@barw}{\linewidth}%
  \par\noindent
  \begin{tikzpicture}[baseline=0pt]
    \useasboundingbox (0,0) rectangle (\igvsc@barw,17pt);
    \clip[rounded corners=3pt] (0,-5pt) rectangle (\igvsc@barw,17pt);
    \fill[vsctabbar] (0,0) rectangle (\igvsc@barw,17pt);
    \node[anchor=south west, inner xsep=7pt, inner ysep=0pt,
          minimum height=17pt, fill=vscbg, text=vsctabtext,
          font=\Lato\fontsize{7.5}{9}\selectfont] (tab) at (0,0)
      {\igvsc@iconmark\hspace{4pt}\lstname\hspace{9pt}%
       {\color{vsctabdim}\fontsize{8}{9}\selectfont\texttimes}};
    \fill[vscaccent] (tab.north west) rectangle ([yshift=-1pt]tab.north east);
  \end{tikzpicture}%
  \par\nointerlineskip}

% La barra compare solo se il listato ha un nome (chiave "name"): listings
% disegna il titolo se \lst@caption non è vuoto, e "title" lo rende non vuoto
% per tutti i listati, quindi lo si svuota quando manca il nome. Va fatto con
% everydisplay, che listings esegue appena prima di disegnare il titolo.
\lstset{title={\lstname}, everydisplay={\igvsc@solosenome}}
\newcommand{\igvsc@solosenome}{%
  \ifx\lst@intname\@empty
    \let\lst@caption\@empty
  \else
    \let\lst@caption\relax
  \fi}
% Al posto della didascalia si disegna la barra delle schede
\AtBeginDocument{\def\lst@maketitle#1{\igvsc@tabbar}}
% Nel nome del file listings sostituisce "-" con un trattino lungo: nella
% scheda deve restare il trattino vero, quello del nome del file.
\def\lst@filenamerpl{_\textunderscore $\textdollar}
\makeatother
